% SSC/ZFS derivation & audit document  (project branch: ssc-zfs)
% Purpose: external-reviewer-auditable derivation of the first-order direct dipolar spin-spin
% contribution to the ZFS D-tensor for MRSF-TDDFT in OpenQP. See CLAUDE.md (esp. §3, §10).
% Mirrored to the Overleaf project at /bighome/vova/coding/SSC/overleaf.
\documentclass[11pt]{article}
\usepackage{amsmath,amssymb,bm,braket,geometry,booktabs,hyperref}
\geometry{margin=1in}

\title{First-Order Direct Spin--Spin Contribution to the Zero-Field-Splitting
$\mathbf{D}$-Tensor for MRSF-TDDFT in OpenQP\\[4pt]\large Derivation and Convention Audit}
\author{OpenQP \texttt{ssc-zfs} project}
\date{\today}

\newcommand{\Pab}{P^{(\alpha-\beta)}}
\newcommand{\rkl}{\dfrac{3\,r_{12,k}\,r_{12,l}-\delta_{kl}\,r_{12}^2}{r_{12}^5}}

\begin{document}
\maketitle

\begin{abstract}
This document re-derives the working equation for the first-order direct dipolar spin--spin (SS)
contribution to the zero-field-splitting (ZFS) $\mathbf{D}$-tensor, resolves the prefactor/sign
convention discrepancy between the source papers, fixes the final prefactor numerically against
the O$_2$ benchmark, and records all assumptions for out-of-band review. It is a living document
updated as the implementation clears its stage gates (L1$\to$L2$\to$L3). \textbf{Status: skeleton.}
\end{abstract}

\tableofcontents

\section{Scope}
First order only; sole physical input is the one-particle spin density $\Pab$ at $M_S=S$.
Z-vector/relaxed densities, gradients of $\mathbf{D}$, and SOC cross terms are out of scope.

\section{Working equation (canonical form)}
McWeeny--Mizuno [Proc.\ R.\ Soc.\ London A \textbf{259}, 554 (1961)], in the form of
Sinnecker \& Neese, J.\ Phys.\ Chem.\ A \textbf{110}, 12267 (2006), Eq.~(9):
\begin{equation}
D^{SS}_{kl} = C \sum_{\mu\nu}\sum_{\kappa\tau}
  \Big\{\Pab_{\mu\nu}\Pab_{\kappa\tau}-\Pab_{\mu\kappa}\Pab_{\nu\tau}\Big\}\,
  \Big\langle \mu\nu \Big| \rkl \Big| \kappa\tau \Big\rangle ,
  \label{eq:working}
\end{equation}
with $k,l\in\{x,y,z\}$; $\mathbf{D}$ symmetric and traceless. The kernel is the Hessian of
$1/r_{12}$ w.r.t.\ the interelectronic vector; the integral is two-electron, ERI-like.

\section{Convention discrepancy and its resolution}\label{sec:conv}
\textbf{The prefactor and sign differ between sources and must not be reconciled by algebra alone}
(see CLAUDE.md~\S3). Recorded verbatim:
\begin{center}
\begin{tabular}{@{}lccc@{}}
\toprule
source & prefactor $C$ & kernel sign & bracket \\
\midrule
Sinnecker--Neese 2006, Eq.~(9) & $+\,g_e^2\alpha^2/[4S(2S-1)]$ & $(3rr-\delta r^2)$ & $\langle\mu\nu|\cdots|\kappa\tau\rangle$ \\
Neese 2007 (JCP 127,164112), Eq.~(46) & $-\,g_e^2\alpha^2/[16S(2S-1)]$ & $(3rr-\delta r^2)$ & $\langle\mu\nu|\cdots|\kappa\tau\rangle$ \\
Neese 2007, Eq.~(1) (operator) & (see text) & $(\delta r^2-3rr)$ & --- \\
\bottomrule
\end{tabular}
\end{center}
Eqs.~(9) and (46) share the kernel sign and density grouping yet differ in $C$ by a factor
$-1/4$; Eq.~(1) flips the kernel sign. \textbf{Resolution (to be filled at gate L2):} adopt the
algebraic \emph{form} of Eq.~\eqref{eq:working}; determine the numerical $C$ (magnitude and sign)
by matching the O$_2$ benchmark. % TODO(L2): insert pinned C, sign, and the match table.

\section{Units}
% TODO: a.u. -> cm^{-1} conversion; role of g_e and alpha; document the exact factor used.

\section{Spin density at $M_S=S$ (Wigner--Eckart)}
% TODO: Pokhilko-Krylov 2019; relate the M_S=S one-particle spin density to the single reduced
% density (highest-spin generator); reuse of soc_mrsf compute_tdm.

\section{Integral evaluation}\label{sec:integral}
\textbf{Route decision (P1.1): Path A.} The two-electron SS integral is the Hessian of $1/r_{12}$
w.r.t.\ the interelectronic vector, so it is computed with the existing ERI machinery rather than a
bespoke Rys primitive. As a distribution,
\begin{equation}
\partial_k\partial_l\frac{1}{r} = \frac{3r_kr_l-\delta_{kl}r^2}{r^5} - \frac{4\pi}{3}\,\delta_{kl}\,\delta^3(\mathbf r),
\label{eq:dist}
\end{equation}
i.e.\ the physical traceless dipolar kernel $T_{kl}=(3r_kr_l-\delta_{kl}r^2)/r^5$ is the
\emph{traceless part} of the bare Hessian; the isotropic contact term does not enter the ZFS
$\mathbf D$-tensor (which is traceless by construction) and is removed.

\paragraph{Closed form for the $(ss|ss)$ quartet.} For unnormalized primitive $s$-Gaussians at
$A,B$ (electron~1, exponents $a,b$) and $C,D$ (electron~2, exponents $c,d$), with the usual
charge-cloud reductions $p=a+b$, $q=c+d$, $P=(aA+bB)/p$, $Q=(cC+dD)/q$, $\rho=pq/(p+q)$,
$R=P-Q$, $T=\rho|R|^2$, $K=e^{-ab|A-B|^2/p}e^{-cd|C-D|^2/q}$, and
$\mathrm{pref}=2\pi^{5/2}/(pq\sqrt{p+q})$, the bare-Hessian integral is
\begin{equation}
H_{kl} = \mathrm{pref}\,K\,\Big[\,4\rho^2 R_k R_l\,F_2(T) - 2\rho\,\delta_{kl}\,F_1(T)\,\Big],
\label{eq:Hclosed}
\end{equation}
($F_n$ the Boys functions), obtained as $H_{kl}=\partial^2 J/\partial d_k\partial d_l|_{d=0}$ from
the operator-displaced Coulomb integral $J(\mathbf d)=\mathrm{pref}\,K\,F_0(\rho|R-\mathbf d|^2)$.
The physical SS integral is its traceless part,
\begin{equation}
S_{kl} = H_{kl} - \tfrac13\,\mathrm{Tr}(H)\,\delta_{kl} = H_{kl} + \tfrac{4\pi}{3}\,\delta_{kl}\,O,
\qquad O=\langle\rho_1|\rho_2\rangle = K\Big(\tfrac{\pi}{p+q}\Big)^{3/2}e^{-T},
\label{eq:Sclosed}
\end{equation}
consistent with Eq.~\eqref{eq:dist}: $\mathrm{Tr}(H)=-4\pi O$, exactly cancelled by the contact term.

\paragraph{Numerical validation (P1.1 prototype, \texttt{tests/ssc\_prototype\_ssss.py}).} The
closed form \eqref{eq:Hclosed}--\eqref{eq:Sclosed} is triangulated against two independent oracles
for a non-degenerate $(ss|ss)$ quartet:
\begin{enumerate}
  \item \textbf{Richardson finite differences} of the Coulomb ERI $J(\mathbf d)$ (using only an
        erf-based $F_0$, independent of the $F_1,F_2$ in \eqref{eq:Hclosed}) $\to$ reproduces $H$.
  \item \textbf{Boys-free Gaussian-transform $t$-quadrature}, $1/r=\frac{2}{\sqrt\pi}\int_0^\infty
        e^{-t^2r^2}dt$, integrated by Gauss--Legendre $\to$ also reproduces $H$ (the contact term is
        captured once the operator is smeared against the clouds; the $t$-quadrature gives the
        \emph{bare} Hessian, \emph{not} the traceless $S$).
\end{enumerate}
Agreement: closed-form $H$ vs FD-of-ERI rel.\ $2.3\times10^{-9}$; closed-form $H$ vs $t$-quadrature
rel.\ $9.9\times10^{-14}$; $\mathrm{Tr}(S)=-3.9\times10^{-16}$; and the contact identity
$\mathrm{Tr}(H)=-4\pi O=-1.904324$ confirmed. \textbf{This validates the closed form against an
analytic oracle only; it is NOT the L1 stage gate} (P1.3), which checks the native Fortran integral
against OpenQP's own ERI engine and is still pending. General angular momenta (P1.2) reuse
\texttt{rys\_deriv.F90}/\texttt{grd2\_rys.F90}.

\section{Assumptions log}
% TODO: every non-obvious choice (reference determinant RO/UNO, mean-field factorization for
% multiconfig MRSF densities, screening, etc.), with justification.
\begin{itemize}
  \item Reference determinant: RO/UNO-type density preferred (UKS overestimates $D^{SS}$;
        ROKS$\approx$UNO -- Neese 2007 \S E). UMRSF exposes a density-source flag.
  \item Integral route: \textbf{Path A} (ERI engine $+$ Hessian-of-$1/r_{12}$ kernel), since the
        SS integral \emph{is} that Hessian (Eq.~\eqref{eq:dist}); no new Rys primitive required
        for validation. The physical kernel is the \emph{traceless part} of the bare Hessian;
        the contact term (Eq.~\eqref{eq:dist}) is dropped as it cannot contribute to a traceless
        $\mathbf D$. Validated numerically in P1.1 (\S\ref{sec:integral}); L1 gate (P1.3) pending.
  \item (add as the project proceeds)
\end{itemize}

\section{Results table}\label{sec:results}
Updated as gates clear (value vs reference vs tolerance, pass/fail).
\begin{center}
\begin{tabular}{@{}lrrrl@{}}
\toprule
system & computed & reference & tol. & gate \\
\midrule
O$_2$\,$^3\Sigma_g^-$ (1.207\,\AA) & --- & $1.44$--$1.6$\,cm$^{-1}$ & few\,\% & L2 \\
benzene & --- & $0.1593$\,cm$^{-1}$ & RMSD$\sim$0.0035 & L3 \\
naphthalene & --- & $0.1004$ & & L3 \\
anthracene & --- & $0.0702$ & & L3 \\
tetracene & --- & $0.0573$ & & L3 \\
\bottomrule
\end{tabular}
\end{center}

\section*{References (key)}
Sinnecker \& Neese, JPCA 110, 12267 (2006);
Neese, JCP 127, 164112 (2007);
Neese, JACS 128, 10213 (2006);
Pokhilko, Epifanovsky \& Krylov, JCP 151, 034106 (2019);
McWeeny \& Mizuno, Proc.\ R.\ Soc.\ A 259, 554 (1961).

\end{document}
